%% PGCT 연구 진행 보고서 (2026-10-22). Overleaf: Compiler = pdfLaTeX (한글 글꼴 오류가 나면 XeLaTeX).
\documentclass[11pt]{article}
\usepackage[a4paper,margin=2.2cm]{geometry}
\usepackage{kotex}
\usepackage{amsmath,amssymb,amsthm}
\usepackage{booktabs}
\usepackage{tabularx}
\usepackage{xltabular}
\usepackage{array}
\usepackage{enumitem}
\usepackage[hidelinks]{hyperref}

\newtheorem{proposition}{명제}
\newtheorem{assumption}{가정}
\newtheorem{definition}{정의}

\newcommand{\JS}{\operatorname{JS}}
\newcommand{\KL}{\operatorname{KL}}
\newcommand{\sg}{\operatorname{sg}}
\newcolumntype{L}{>{\raggedright\arraybackslash}X}
\setlist{itemsep=2pt,topsep=4pt}
\renewcommand{\arraystretch}{1.15}

\title{Probe-Gated Critic Transfer 연구 진행 보고서}
\author{이환률}
\date{2026년 10월 22일}

\begin{document}
\maketitle

\begin{abstract}
PGCT는 팀 소속을 모르는 agent가 공통 probe에 대한 이웃의 정책 응답을 비교해 팀원을 추정하고, 추정된 팀원의 critic만 증류하는 방법이다.
지금까지의 실험에서 정책 probe는 어느 환경에서도 팀을 안정적으로 가르지 못했다.
또 환경의 pickup 감점 때문에 학습 초반에 선반 집기가 사라져, 방법 간 차이가 ``누가 먼저 재발견하나''라는 운에 묻혔다.
감점을 .005로 낮추어 이 정체를 풀고 있다.
팀 보상 환경에서는 ``남의 배송이 내 보상을 바꾸는가''를 묻는 보상 결합 probe가 팀을 완벽히 갈랐다(AUROC 1.00).
이에 따라 연구 목표(잠재 팀 발견 + 팀 내 critic 전이)는 유지한다.
대신 기여를 ``어떤 조건에서 국소 신호로 팀을 찾을 수 있는가''로 다시 잡고, 같은 기반 위의 2$\times$2 실험으로 검증한다.
\end{abstract}

\tableofcontents

%% ------------------------------------------------------------------
\section{연구 개요}

\paragraph{문제.} 같은 창고를 두 회사의 로봇이 함께 쓴다고 하자. 각 로봇은 자기 팀이 요청한 선반을 자기 팀 배송 지점으로 옮긴다.
협력 MARL에서는 목표가 같은 agent끼리 파라미터·경험·가치 추정을 나누면 표본 효율이 좋아진다.
그러나 다른 팀의 critic을 배우면 그 critic은 다른 목표의 가치 쪽으로 편향된다.
팀원이 누구인지는 항상 주어지지 않는다. 공통 식별 체계가 없을 수 있고, 배치마다 팀이 바뀔 수 있다.

\paragraph{목표.} 국소 정보만으로 잠재 팀원을 찾고, 그 팀원과만 critic을 공유해 학습을 개선한다.
중앙 서버와 전역 팀 라벨이 없는 분산 환경을 가정한다. 유용한 critic을 찾는 것은 목적이 아니라 검증 과정이다(D41, D54).

\paragraph{세 질문.}
\begin{itemize}
  \item Q1: agent는 국소 신호만으로 이웃 가운데 팀원을 식별할 수 있는가?
  \item Q2: 팀원의 critic은 실제로 학습에 도움이 되는가?
  \item Q3: Q1의 올바른 답이 학습 향상으로 이어지는가?
\end{itemize}

\paragraph{일정.} 목표 학회는 IJCAI 2027이다. 초록 마감은 2027년 1월 4일, 본문 마감은 1월 11일(AoE)이다.

%% ------------------------------------------------------------------
\section{PGCT 방법}

\subsection{문제 설정}
agent $i \in \{1,\dots,N\}$로 이루어진 부분 관측 확률 게임을 생각한다. 각 agent는 팀 $z_i \in \{1,\dots,K\}$에 속한다.
팀은 어떤 선반이 요청되는지와 어디로 배송하는지를 정한다.
agent의 관측은 자기 팀이 필요로 하는 정보를 담지만, 다른 agent의 팀 $z_j$는 담지 않는다.
통신 범위 안의 agent가 이웃 $\mathcal{N}_i(t)$이다.
이웃끼리는 probe 입력에 대한 네트워크 출력만 교환한다. 관측·보상·파라미터·팀 라벨은 교환하지 않는다.
각 agent는 순환 actor $\pi_i$와 critic $V_i$를 따로 두고, 자기 경험으로 PPO 학습한다.

\begin{definition}[부분 팀 관계]
$\mathcal{T}_i(t) = \{ j \in \mathcal{N}_i(t) : z_j = z_i \}$. 모든 $i$에 대해 추정 $\hat{\mathcal{T}}_i(t) = \mathcal{T}_i(t)$이면, 그 방법은 시점 $t$에서 부분 팀 관계를 복원한다고 한다.
\end{definition}

\subsection{Probe 집합}
probe는 작고 통제된 상황 하나를 나타내는 관측 하나다. 예를 들면 ``자기 팀 요청 선반을 운반 중이고 배송 지점은 시야 밖''인 상황이다.
모든 probe는 학습 전에 지도 배치로부터 정해진 규칙으로 만든다. 그래서 모든 agent가 통신 없이 같은 집합을 갖는다.
probe는 순환 상태를 0으로 두고 action mask 없이 평가한다. 질의 한 번이 forward 한 번이고, 학습 궤적에는 영향을 주지 않는다.

\begin{assumption}[probe 비교 가능성]
모든 probe는 모든 agent에게 유효한 관측이며, 관측하는 agent 자신의 목표에 대한 상황을 모두에게 같은 방식으로 나타낸다.
\end{assumption}

probe는 두 역할을 한다. gate 집합 $\mathcal{Q}_{\mathrm{gate}}$에는 팀에 따라 최적 행동이 달라지는 상황이 필요하다.
증류 집합 $\mathcal{Q}_{\mathrm{dist}}$에는 과제의 모든 단계를 덮는 넓은 범위가 필요하다.
공유 설정은 $\mathcal{Q}_{\mathrm{gate}} = \mathcal{Q}_{\mathrm{dist}}$(48개, 서로 다른 입력 26개)이고, 분리 설정은 gate 48개와 증류 64개다.

\subsection{응답 비교}
$p_i^m = \pi_i(\cdot \mid q_m, h = 0)$이라 하자. 이웃은 $P$ update마다 $\{p_i^m\}$을 교환하고 다음을 계산한다.
\[
  D_{ij} = \frac{1}{M}\sum_{m=1}^{M} \JS(p_i^m, p_j^m), \qquad
  \JS(p,q) = \tfrac12 \KL(p \,\|\, \bar p) + \tfrac12 \KL(q \,\|\, \bar p), \quad \bar p = \tfrac{p+q}{2}.
\]
JS는 대칭이고, 한쪽 확률이 0이어도 유한하며, $[0, \ln 2]$에 갇힌다. 정책이 거의 결정적이면 $D_{ij}$는 응답이 다른 probe의 비율에 $\ln 2$를 곱한 값에 가깝다.

\subsection{Gate}
한 쌍의 첫 측정으로 $\bar D_{ij} \leftarrow D_{ij}$를 두고, 이후에는 $\bar D_{ij} \leftarrow (1-\beta)\bar D_{ij} + \beta D_{ij}$로 평활한다.
측정은 두 agent가 이웃일 때만 한다. 측정 횟수를 $n_{ij}$에 센다. warm-up $U_0$ update가 지나면 다음을 팀원으로 본다.
\[
  \hat{\mathcal{T}}_i(t) = \{ j \in \mathcal{N}_i(t) : n_{ij} \ge n_{\min},\ \bar D_{ij} < \tau \}.
\]

\subsection{Critic 증류}
\[
  \mathcal{L}_i = \mathcal{L}^{\mathrm{PPO}}_i + \lambda \sum_{j \in \hat{\mathcal{T}}_i(t)} \alpha_{ij}\,
  \frac{1}{M_d}\sum_{m=1}^{M_d} \Big( V_i(q^d_m) - \sg\big[V_j(q^d_m)\big] \Big)^2, \qquad \alpha_{ij} = \frac{1}{|\hat{\mathcal{T}}_i(t)|}.
\]
$V_j$는 update 전에 계산한 고정 목표다. $\hat{\mathcal{T}}_i(t)$가 비면 peer 항이 사라지고 독립 PPO와 같아진다.
actor 대신 critic을 전이하는 이유는 두 가지다. 상황의 가치는 팀 목표의 속성이지만 행동 선택은 역할에도 달려 있다. 또 actor는 advantage $\hat A = \hat R - V_i$를 통해서만 간접적으로 영향을 받는다.

\subsection{알고리즘}
agent $i$는 다음을 반복한다.
\begin{enumerate}
  \item 초기화: $\pi_i, V_i$를 만들고, 모든 $j$에 대해 $\bar D_{ij}$는 미정, $n_{ij} \leftarrow 0$.
  \item update $u = 1, 2, \dots$마다
  \begin{enumerate}
    \item $\pi_i$로 rollout을 모으고 자기 보상으로 PPO 목표를 계산한다.
    \item $u \bmod P = 0$이면, 이웃 $j$마다 $\mathcal{Q}_{\mathrm{gate}}$의 응답을 교환해 $D_{ij}$를 구하고 $\bar D_{ij}$와 $n_{ij}$를 갱신한다.
    \item $u > U_0$이면 $\hat{\mathcal{T}}_i$를 계산하고, 아니면 $\hat{\mathcal{T}}_i = \emptyset$으로 둔다.
    \item $j \in \hat{\mathcal{T}}_i$로부터 update 전 값 $V_j(q^d_m)$을 받는다.
    \item $\mathcal{L}_i$로 $\pi_i, V_i$를 갱신한다.
  \end{enumerate}
\end{enumerate}
설정값은 $P = 5$, $\beta = 0.35$, $\tau = 0.12$, $U_0 = 100$, $n_{\min} = 8$, $\lambda = 0.01$이다.

\subsection{비교 방법}
모든 방법은 같은 순환 PPO 학습기(semantic CNN 인코더, 128 unit GRU)와 같은 하이퍼파라미터를 쓴다. seed가 같으면 초기 가중치도 비트 단위로 같다.
\begin{itemize}
  \item IPPO: 전이 없음.
  \item All-peer: 모든 이웃을 donor로 쓴다.
  \item Oracle-team: 이웃 중 실제 팀원만 쓴다. 사실상 Du et al.\ \cite{du2025}의 목표 인지 공유(critic만)에 해당한다.
  \item Random-gated: update마다 PGCT가 실제로 만든 연결 수만큼 이웃을 무작위로 고른다. 전이량을 통제하기 위한 대조군이다(E35).
  \item PGCT, PGCT-split: 공유 설정과 분리 설정.
\end{itemize}

%% ------------------------------------------------------------------
\section{이론}

\paragraph{분리 마진.} 측정 $n$에서 이웃 쌍에 대해 $\delta_{\mathrm{in}}(n) = \max_{z_i = z_j} D_{ij}(n)$, $\delta_{\mathrm{out}}(n) = \min_{z_i \ne z_j} D_{ij}(n)$, $\Delta(n) = \delta_{\mathrm{out}}(n) - \delta_{\mathrm{in}}(n)$로 둔다.

\begin{proposition}[정확한 부분 복원]
$\Delta(n) > 0$이고, 모든 이웃 쌍에서 $|\bar D_{ij}(n) - D_{ij}(n)| \le \epsilon$, $2\epsilon < \Delta(n)$이며, warm-up과 최소 측정 횟수가 충족되고 $\delta_{\mathrm{in}} + \epsilon < \tau < \delta_{\mathrm{out}} - \epsilon$이면, 모든 $i$에 대해 $\hat{\mathcal{T}}_i(n) = \mathcal{T}_i(n)$이다.
\end{proposition}

\begin{proposition}[추적 오차]
연속된 측정의 차이가 $|D_{ij}(n) - D_{ij}(n-1)| \le \rho$이면, 첫 측정으로 시작한 이동평균은 $|\bar D_{ij}(n) - D_{ij}(n)| \le \frac{1-\beta}{\beta}\rho$를 만족한다. 따라서 $2\frac{1-\beta}{\beta}\rho < \Delta(n)$이면 명제 1이 적용된다. $\beta = 0.35$이면 $(1-\beta)/\beta \approx 1.86$이다.
\end{proposition}

\begin{proposition}[같은 결정 문제를 푸는 팀]
다른 agent의 정책을 고정했을 때 서로 다른 팀 agent의 결정 문제가 같으면, 기억 없는 정책 가운데 한 팀에 최적인 사상은 다른 팀에도 최적이다. 최적 사상이 유일하면 모든 쌍에서 $D_{ij} = 0$, $\Delta = 0$이 되어 어떤 임계값도 두 팀을 가르지 못한다.
\end{proposition}

\paragraph{확장: 식별 가능성의 두 원천.} 명제 3의 논리는 비교하는 양이 관측자 기준 $(o, a)$의 함수이기만 하면 그대로 성립한다. 정책이든 critic이든 보상 예측기 $\hat r_i(o,a)$든 마찬가지다. 팀을 가를 수 있는 원천은 두 가지로 정리된다.
\begin{itemize}
  \item (가) 목표의 기하: 공통 좌표(절대 위치)에서 팀마다 목표가 다르면 행동·가치 probe로 가를 수 있다.
  \item (나) 보상의 결합: 남의 사건이 내 보상을 바꾸면(팀 보상) 보상 결합 probe로 가를 수 있다.
\end{itemize}
둘 다 없으면 어떤 국소 신호로도 팀을 가를 수 없다. 이 구분이 이후 실험 설계의 기준이다.
행동 유사성은 같은 팀의 필요조건도 충분조건도 아니다. 관측이 ``내 기지, 내 자원'' 기준이면 다른 팀의 일꾼도 똑같이 행동한다. 같은 팀이라도 역할을 나누면 다르게 행동한다.

\subsection{보상 결합 probe의 정의}
PGCT에서 바뀌는 것은 gate 신호 하나다. 정책 응답의 JS 거리 대신 공통 사건에 대한 ``내 보상 반응''의 거리를 쓰고, critic 증류 손실은 그대로 둔다. 아래는 E49·E50에서 구현한 형태다.

\paragraph{공통 사건 집합.} 지도에서 정해진 규칙으로 사건 목록 $\mathcal{E} = \{e_1, \dots, e_M\}$을 만든다. 모든 agent가 통신 없이 같은 목록을 갖고, 팀 라벨은 쓰지 않는다. 지금은 goal 칸 $g$마다 사건 $e_g$(``칸 $g$에서 배송이 일어남'') 하나씩, $M = 2$다.

\paragraph{agent별 통계.} 사건 $o_i(e,t)$를 ``agent $i$가 시점 $t$에 사건 $e$를 봤다''로 정의한다. 일으킨 agent가 $i$가 아니고, 사건 위치가 $i$의 시야(센서 범위 $R_s = 5$) 안이며, 그 step에 $i$ 자신은 배송·반납을 하지 않았다는 뜻이다.
\[
  S_i(e) = \sum_t \mathbb{1}[o_i(e,t)]\, r_i(t), \qquad
  C_i(e) = \sum_t \mathbb{1}[o_i(e,t)], \qquad
  \rho_i(e) = \frac{S_i(e)}{C_i(e)} - b_i \quad (C_i(e) \ge K).
\]
$b_i$는 아무 사건도 보지 않았고 자기 사건도 없던 step의 평균 보상이다. $\rho_i(e)$는 ``$e$가 일어났을 때 내 보상이 평소보다 얼마나 달라지나''의 추정치이며, 이것이 probe 응답이다. 시뮬레이션도 신경망 학습도 없고, 자기 경험의 평균일 뿐이다.

\paragraph{비교와 판정.} 두 agent가 통신 범위 안에서 만나면 $\rho$ 벡터와 횟수 $C$를 교환한다. 보상 자체는 교환하지 않는다.
\[
  \mathcal{M}_{ij} = \{e : C_i(e) \ge K,\ C_j(e) \ge K\}, \qquad
  D^r_{ij} = \frac{1}{|\mathcal{M}_{ij}|} \sum_{e \in \mathcal{M}_{ij}} \big|\rho_i(e) - \rho_j(e)\big|,
\]
\[
  \hat{\mathcal{T}}_i = \{\, j \in \mathcal{N}_i(t) : \mathcal{M}_{ij} \ne \emptyset,\ D^r_{ij} < \theta \,\}, \qquad \theta = R/2.
\]
$R$은 배송 보상(여기서는 1)이다. 판정은 다음 만남까지 유지하고, 판정 전에는 전이하지 않는다. 이후 증류는 2.5절과 같다.

\paragraph{분리 조건 (스케치, 정식 증명 아님).} 팀 보상에서는 goal $g$에 배송이 일어나는 순간 $g$를 가진 팀 전원이 $R$을 받으므로 $\rho_i(e_g) \approx R \cdot \mathbb{1}[g \in G_{z_i}]$이다. 추정 오차를 $|\hat\rho - \rho| \le \epsilon$이라 하면
\[
  z_i = z_j \Rightarrow |\rho_i(e) - \rho_j(e)| \le 2\epsilon, \qquad
  z_i \ne z_j \Rightarrow |\rho_i(e) - \rho_j(e)| \ge R - 2\epsilon
\]
가 goal 칸이 팀마다 하나인 경우 모든 $e$에서 성립한다. 따라서 $\epsilon < R/4$이면 $\theta = R/2$가 두 팀을 정확히 가른다. 명제 1과 같은 구조로, 분리 마진이 $R - 4\epsilon$인 셈이다. $\epsilon$은 shaping 잡음과, 같은 step에 시야 밖에서 일어난 다른 배송에서 생기며 $K$가 커질수록 줄어든다.

\paragraph{실패 조건.} 필요조건은 어떤 사건 $e$에 대해 다른 팀끼리 $\mathbb{E}[r_i \mid e] \ne \mathbb{E}[r_j \mid e]$인 것이다(조건 (나)).
\begin{itemize}
  \item 개인 보상: 남의 배송이 내 보상을 바꾸지 않아 모든 $\rho \approx 0$이다(E42에서 쌍별 $|r| \le 0.015$).
  \item 두 팀이 goal 칸을 공유: 그 칸의 배송은 선반 주인에 따라 보상받는 팀이 달라, 두 팀 모두 $\rho \approx 0.5$ 근처가 된다.
  \item 보상 지연이나 큰 잡음: $\epsilon$이 $R/4$를 넘을 수 있다.
\end{itemize}

\paragraph{변형.} 직접 증거형 $c_{ij} = \mathbb{E}[r_i \mid j\text{의 사건을 봄}] - b_i$는 goal을 공유해도 쓸 수 있지만, 다른 팀 선반을 대신 배송하는 경우에 틀린다(전역 관측 상한 0.83, E49). 사건이 복잡하면 표 대신 회귀 모델 $f_i(\phi(e)) \approx \mathbb{E}[r_i \mid e]$를 학습해 같은 거리로 비교한다. 이 일반화형은 아직 시험하지 않았다.

%% ------------------------------------------------------------------
\section{실험 경과}

\subsection{환경}
다중 로봇 창고(RWARE) \cite{papoudakis2021}의 두 팀 변형이다. 지도는 11$\times$10 격자에 선반 32개가 있다.
agent는 11$\times$11 시야를 13개 채널로 보고, 절대 위치·방향·적재 상태를 포함한 20개 정보를 함께 받는다.
채널은 팀 기준이다. 자기 팀 요청 선반과 자기 배송 지점이 전용 채널에 표시된다.
보상은 배송 +1, 반납 +1, 요청 선반 적재 +0.15와 작은 shaping 항이다. episode는 500 step이다.
배치는 E1(배송 지점 인접, (4,10)/(5,10))과 E1-F(양쪽 모서리, (0,10)/(9,10))다.
평가는 결정적 정책으로 20 episode를 돌려 배송 수를 센다.

\subsection{E1 (4 agent, 개인 보상, seed 1, 6.4M)}
\begin{center}
\begin{tabular}{lccccc}
\toprule
방법 & 시작 $\ge 1$ (M) & 시작 $\ge 20$ (M) & 평균 0--6.4M & 마지막 5회 & 학습 agent \\
\midrule
IPPO & 2.05 & 5.07 & 10.4 & 23.2 $\pm$ 1.1 & 2/4 \\
All-peer & 3.99 & -- & 3.7 & 19.4 $\pm$ 1.6 & 4/4 \\
PGCT & 1.79 & 4.15 & 11.7 & 27.3 $\pm$ 1.4 & 3/4 \\
Oracle-team & 1.64 & 4.66 & 13.2 & 32.2 $\pm$ 1.7 & 3/4 \\
\bottomrule
\end{tabular}
\end{center}
PGCT가 IPPO보다 빨라 보이지만, gate는 팀원을 고르지 못했다. E1 checkpoint에서 같은 팀 probe 거리는 0.333, 다른 팀은 0.276으로 오히려 반대였다.
같은 seed에서도 첫 발견 시점이 0.2--2.1M으로 흔들려, seed 하나의 성능 차이는 발견 운이 증폭된 값일 수 있다(D45).

\subsection{Probe 점검과 보상 probe (E29--E39)}
\begin{itemize}
  \item 정책 JS probe: agent 자신의 정책 표류가 팀 차이만큼 컸다. 특정 계보에서만 나오는 신호는 독립 계보로 일반화되지 않았다(E29--E34).
  \item 행동 궤적(시퀀스) probe: 길이 1 대비 이점이 없었다(E37, E38).
  \item 보상 probe $\hat r(o,a)$: 인위적 probe 상태에서는 팀별 실제 보상이 다른 query가 2\%뿐이라 실패했다(E36). 실제 방문 상태에서만 신호가 남았는데, 그 신호도 주로 각 agent가 자기 팀 보상을 배운 결과였다(E39).
\end{itemize}

\subsection{E1-F 중단과 환경 재검토 (E26, E28, E40--E43)}
E1-F 주 실험 5개는 10월 19일 Claude 앱 업데이트로 1.536M에서 멈췄다. 그때까지 평가 배송은 거의 0이었고 gate는 같은 팀을 고르지 않았다.
E1-F는 개인 보상이라 팀원 critic도 같은 값을 추정하지 않는다. 그래서 그대로 재시작하지 않고 팀 보상 후보 A(4 agent)와 B(6 agent, 통신 범위 3)를 만들었다.
IPPO 2M pilot(E43)에서 A seed 1은 배송이 없었고, A seed 2는 1.89M에 한 agent만 배웠다. B seed 1은 1.13M에 한 agent만 배웠고, B seed 2는 배송이 없었다.

\subsection{Pickup 정체의 원인 (E44)}
학습 초기 정책은 거의 무작위다. 무작위 정책에서 행동별 평균 즉시 보상을 재면 선반 집기(TOGGLE)가 가장 손해다.
\begin{center}
\begin{tabular}{lcccc}
\toprule
행동 & FORWARD & LEFT/RIGHT & NOOP & TOGGLE \\
\midrule
평균 즉시 보상 & $-0.0002$ & $-0.0022$ & $-0.0053$ & $-0.0183$ \\
\bottomrule
\end{tabular}
\end{center}
TOGGLE의 45\%는 요청되지 않은 선반 집기($-0.02$), 48\%는 조기 내려놓기($-0.02$)다. 자기 팀 요청 선반을 집는 경우는 3.6\%뿐이다.
그래서 모든 run에서 0.1--0.2M 안에 선반 집기가 사라지고, 우연히 다시 발견한 agent 하나만 배운다. 이것이 ``누가 먼저 발견하나'' 잡음의 정체다.

\subsection{감점 조정 (E46--E48)}
\begin{itemize}
  \item 감점 0(nopen, E46): TOGGLE과 요청 선반 pickup은 1M 내내 유지되어 원인이 인과적으로 확인됐다. 대신 요청되지 않은 선반 pickup이 50--150/episode로 폭증했고, 1M 안에 평가 배송 $\ge 1$에 도달한 run이 없었다.
  \item 감점 크기 sweep(E47, 학습 없음): 무작위 정책에서 TOGGLE 즉시 보상은 대략 $0.0002 - 0.93p$다. NOOP보다 나아지는 경계는 $p \approx .0056$--$.0059$다.
  \item 감점 .005(pen005, E48, 진행 중): 0.1--0.2M 구간 TOGGLE/episode가 87--106으로, 기존 .02(3--9)와 0(약 270) 사이에 있다. 10월 22일 14:28 기준 약 0.27M이다.
\end{itemize}

\subsection{보상 결합 probe 대 정책 probe (E49, B nopen)}
보상 결합 probe는 지도에서 정한 공통 사건 ``goal 칸 $g$에서 배송''을 쓴다. 각 agent는 시야 안에서 남이 그 사건을 일으킨 step마다 받은 자기 보상을 평균한다(기저 보상을 뺀 값). 이웃은 이 응답만 비교한다. 실제로 goal에 넣어 보는 시뮬레이션도, 신경망 예측기도 아니다. 정의는 3.1절에 있다.
\begin{center}
\begin{tabular}{lcccc}
\toprule
팀 AUROC & s1 0.77M & s1 1M & s2 0.77M & s2 1M \\
\midrule
보상 결합 probe & 1.00 & 1.00 & 1.00 & 1.00 \\
정책 probe (canonical 48) & 0.78 & 0.67 & 0.74 & 0.59 \\
정책 probe (v2 gate 48) & 0.56 & 0.76 & 0.76 & 0.41 \\
정책 probe (방문 상태 256) & 0.78 & 0.44 & 0.74 & 0.56 \\
\bottomrule
\end{tabular}
\end{center}
1.00은 방법이 뛰어나서가 아니다. 팀 보상이 배송 순간 팀 전원에게 잡음 없이 들어가고, goal 칸 하나가 정확히 한 팀에 속하기 때문이다. 이 환경에서 이 신호는 사실상 팀의 정의다.
또 배송한 agent와 보상받는 팀이 다를 수 있다(다른 팀 선반을 대신 배송). 두 팀이 같은 goal을 공유하면 위치 기준 probe도 갈리지 않는다.

\subsection{식별 시점 (E50)}
episode를 순서대로 돌리며 온라인으로 판정했다. 두 agent가 통신 범위 안에서 만나고 같은 사건을 $K$번 이상 봤을 때만 판정하며, 판정 전에는 전이하지 않는다.
\begin{center}
\begin{tabular}{lcccc}
\toprule
15쌍 모두 정답인 env step & s1, $K{=}1$ & s1, $K{=}3$ & s2, $K{=}1$ & s2, $K{=}3$ \\
\midrule
무작위 정책 & 15,283 & 미달 (10/15) & 미달 (10/15) & 미달 (6/15) \\
0.26M & 36,472 & 미달 (11/15) & 10,987 & 19,289 \\
0.51M & 12,944 & 32,785 & 6,106 & 33,749 \\
0.77M & 19,551 & 35,624 & 9,922 & 미달 (13/15) \\
1M & 11,034 & 22,200 & 20,501 & 미달 (14/15) \\
\bottomrule
\end{tabular}
\end{center}
``미달''은 4만 step 안에 다 찾지 못한 경우다. 오판정은 20개 조건 중 1회($K{=}1$)뿐이었다. $K{=}1$이면 대체로 1--4만 step, 즉 1M 학습의 1--4\% 안에 팀을 다 찾는다.
다만 nopen은 TOGGLE 남발로 우연한 배송이 잦다. 감점이 있는 환경에서는 더 늦어질 수 있어, pen005에서 다시 잰다.

%% ------------------------------------------------------------------
\section{받은 질의와 답}

\begin{xltabular}{\textwidth}{>{\raggedright\arraybackslash}p{4.2cm} L >{\raggedright\arraybackslash}p{2.6cm}}
\toprule
질의 & 답 & 근거 \\
\midrule
\endhead
\textbf{독창성}: 같은 목표의 이웃과만 공유하는 방법은 이미 있지 않나? & 있다. Du et al.\ \cite{du2025}은 목표를 명시적으로 교환하고 같은 목표의 이웃과 actor·critic 가중치를 평균한다. PGCT의 차이는 목표·라벨을 교환하지 않고 공통 probe의 출력만 비교한다는 점이다. & 원문 확인 \\
\textbf{독창성}: 그럼 gate만 바꾼 것 아닌가? & 그렇게 읽힐 수 있다. 그래서 기여를 국소 신호로 잠재 팀을 찾을 수 있는 조건과, 조건에 맞는 probe 두 종류로 다시 잡는다. & D71 \\
\textbf{독창성}: ``도움이 되는 상대 찾기''로 바꾸면? & 안전하지 않다. Peer Learning \cite{derstroff2024}이 유용성 기반 선택을 이미 다룬다. 연구 목적(팀 발견)은 유지한다. & D41, D54 \\
\textbf{알고리즘}: 행동이 비슷하면 같은 팀이라는 보장이 있나? & 없다. 팀 기준 관측에서는 다른 팀도 똑같이 행동하고, 역할을 나누면 같은 팀도 다르게 행동한다. 전제를 조건 (가)로 낮춘다. & 명제 3, E1 \\
\textbf{알고리즘}: 보상 기반 probe가 낫지 않나? & 방향은 맞지만 형태가 중요하다. $\hat r(o,a)$는 행동 probe와 같은 함정에 빠진다. 보상 결합 probe는 팀 보상에서 1.00이지만, 개인 보상에서는 결합이 거의 0이다. & E36--E39, E42, E49 \\
\textbf{알고리즘}: 고정 $\tau$와 $h{=}0$ probe는 믿을 만한가? & 약하다. 정책 표류가 팀 신호만큼 크고, probe는 학습 분포 밖이다. 이웃 간 상대 순위 판정을 검증 후보로 둔다. & E29--E34, E49 \\
\textbf{알고리즘}: 왜 critic만 전이하나? & 가치는 팀 목표의 속성이고 행동은 역할에도 달려 있다. 다만 개인 보상에서는 팀원 critic도 같은 값을 추정하지 않으므로, 전이 효과는 팀 보상 조건에서 본다. & 2.5절 \\
\textbf{알고리즘}: 가정이 너무 많다 & 맞다. 전제는 조건으로 바꾸고 명제 1·2는 부록으로 옮긴다. 남는 핵심은 가정 1과 공통 probe 생성 규칙이다. & D71 \\
\textbf{실험}: seed 하나로는 부족하다 & 맞다. seed 3개로 시작해 5개로 늘리고 IQM과 신뢰구간을 보고한다 \cite{agarwal2021}. & D45 \\
\textbf{실험}: All-peer보다 나은 게 상대 선택 때문인가, 전이를 덜 해서인가? & 구분할 수 없었다. 전이량을 맞춘 Random-gated를 구현했다. & E35 \\
\textbf{실험}: 목표가 다른 조건(E1-F) 결과가 없다 & 1.536M에서 중단됐다. 개인 보상과 pickup 정체 문제가 있어 환경부터 고쳤다. & E26, D62 \\
\textbf{실험}: 조건 간 비교가 공정한가? & 아니었다. E1·E1-F와 B는 agent 수·통신 범위·감점이 달랐다. 같은 기반의 2$\times$2로 다시 짰다. & D74, E51 \\
\textbf{실험}: baseline이 좁다 & 일반 성능 우위를 주장하지 않는다면 필수는 IPPO, All-peer, Oracle-team, Random-gated다. MAPPO·SePS는 구현되어 smoke만 했다. & 8절 \\
\bottomrule
\end{xltabular}

%% ------------------------------------------------------------------
\section{문제점 정리}

\begin{xltabular}{\textwidth}{>{\raggedright\arraybackslash}p{2.2cm} L L}
\toprule
영역 & 문제 & 영향 \\
\midrule
\endhead
환경 & 감점 .02에서 TOGGLE이 초기 탐색의 최악 행동이라 0.1--0.2M에 선반 집기가 사라진다 & 한 agent만 배우고, 결과가 발견 운에 좌우된다 \\
환경 & 감점 0이면 TOGGLE을 남발한다 & 본 실험 환경으로 부적합하다 \\
지표 & seed 하나의 성능 비교 & E1 결과를 근거로 쓸 수 없다 \\
gate & 신호 반전(같은 팀 0.333 > 다른 팀 0.276), 2M 이후 팀원 간 전이 소멸 & 정책 probe가 어느 환경에서도 팀을 안정적으로 가르지 못했다 \\
gate & $h{=}0$·mask 없는 probe, 중복 입력, 가정 1 위반 probe & 분포 밖 응답을 비교한다 \\
독창성 & Du et al.과 가깝다 & gate 메커니즘만으로는 기여가 약하다 \\
가정 & 불필요한 가정이 많다 & 공격 리뷰의 주된 표적이 된다 \\
실험 설계 & 조건마다 기반 설정이 다르다 & 식별 성패의 원인을 가를 수 없다 \\
초판 & E1-F 상태 오기, 깨진 참조, 형식 오류 & 발표본이 현재와 다르다 \\
\bottomrule
\end{xltabular}

\paragraph{E1-F에 대한 주의.} E1-F에서 정책 probe가 실패했다고 해서 조건 (가)가 반박된 것은 아니다. 1.5M까지 평가 배송이 거의 0이어서 정책이 아직 자기 goal로 가는 법을 배우지 못한 상태였다. 목표가 갈라져 있으면서 여러 agent가 배송을 배운 run은 아직 없다.

%% ------------------------------------------------------------------
\section{수정 방향}

\paragraph{연구 구도.} 중심은 팀 발견으로 둔다. 기여는 gate 메커니즘이 아니라, 어떤 팀 구조에서 어떤 국소 신호로 팀을 찾을 수 있는가라는 조건과 조건에 맞춘 probe다.

\paragraph{식별 가능성 2$\times$2.} 기반을 6 agent(3+3), 센서 5, 통신 범위 3, 두 감점 .005로 고정하고, goal 배치와 보상 방식만 바꾼다. 설정 대조에서 다른 값은 \texttt{team\_reward\_mode}와 \texttt{team\_goal\_positions}뿐이었다.
\begin{center}
\begin{tabularx}{\textwidth}{>{\raggedright\arraybackslash}p{2.6cm} L L}
\toprule
goal 배치 & 개인 보상 & 팀 보상 \\
\midrule
인접 (4,10)/(5,10) & 둘 다 없음. 어떤 신호도 못 가를 것으로 예상 (\texttt{b6\_adjacent\_indiv\_pen005}) & (나)만. 정책 probe 실패, 보상 결합 성공 예상 (\texttt{b6\_adjacent\_team\_pen005}) \\
모서리 (0,10)/(9,10) & (가)만. 정책 probe만 쓸 수 있는 식별의 승부처 (\texttt{b6\_corner\_indiv\_pen005}) & 둘 다 있음. 전이 효과를 보기 가장 좋은 칸 (\texttt{b\_e1f6\_comm3\_pen005}) \\
\bottomrule
\end{tabularx}
\end{center}

\paragraph{구체적 변경.}
\begin{itemize}
  \item ``같은 목표면 비슷하게 행동한다''를 전제에서 조건 (가)로 바꾼다.
  \item 명제 1·2는 부록으로 옮기고, 명제 3은 정책·가치·보상 예측 모두에 해당하도록 넓힌다.
  \item 고정 $\tau$ 대신 이웃 간 상대 순위 판정을 검증한다.
  \item 보상 결합 gate는 학습용으로 구현하지 않는다. 팀 보상 칸에서는 학습의 1--4\% 안에 Oracle-team과 같아지고, 개인 보상 칸에서는 작동하지 않는다. 두 결과 모두 오프라인 측정으로 보인다.
  \item 관련 연구에 Du et al.\ \cite{du2025}, Peer Learning \cite{derstroff2024}, IDRL \cite{idrl} 등을 넣는다.
  \item 전이 효과(Q2)는 팀 보상 칸에서 본다. 팀원이 같은 return을 추정하는 조건이기 때문이다. 개인 보상 칸에서는 팀원끼리 요청을 두고 경쟁해, Oracle-team이 All-peer보다 낫다는 보장이 없다.
\end{itemize}

%% ------------------------------------------------------------------
\section{실험 계획과 일정}

각 단계는 앞 단계의 판정을 보고 사용자가 직접 시작한다. 모든 학습은 Claude 앱 밖 PowerShell에서 돌린다.
\begin{enumerate}
  \item \textbf{pen005 판정} (10월 22일 자정 전후 종료, 10월 23일 확인). 판정 스크립트 \texttt{latent\_team/pilot\_verdict.py}가 아래 기준으로 판정하고 보고서를 쓴다.
  \item \textbf{2$\times$2 훑기} (10월 22일 14:01 시작, 10월 23일 밤 종료 예상). 새 칸 3개 $\times$ seed 2, IPPO 1M을 대기열(\texttt{latent\_team/run\_queue.ps1}, 동시 4개)로 돌린다. 끝나면 칸마다 정책 probe AUROC, 보상 결합 AUROC, 식별 시점을 학습 없이 잰다.
  \item \textbf{메인 비교 1단계} (약 5일). 모서리·팀 칸과 모서리·개인 칸에서 IPPO, All-peer, Oracle-team, PGCT-정책 $\times$ seed 3. 길이는 훑기에서 배송 시작 시점을 보고 정한다(지금 예상 3M).
  \item \textbf{메인 비교 2단계} (약 2일). Random-gated $\times$ seed 3. 끝난 PGCT run의 전이량에 맞춰야 하므로 1단계 뒤에 돌린다.
  \item \textbf{11월}: seed 4·5 추가, 결정적인 칸의 길이 연장, IQM과 신뢰구간.
  \item \textbf{12월}: 원고 작성과 보완 실험.
\end{enumerate}

\begin{center}
\begin{tabularx}{\textwidth}{>{\raggedright\arraybackslash}p{4.2cm} L}
\toprule
pilot 판정 기준 (두 seed, 1M 기준) & 통과 조건 \\
\midrule
1. TOGGLE 유지 & 0.2M 이후 모든 0.1M 구간에서 TOGGLE/episode $\ge 5$ \\
2. 요청 선반 pickup & 0.2M 이후 평균 $\ge 1$/episode \\
3. 배송 & 한 seed라도 평가 배송 $\ge 1$이 3회 연속, 또는 두 seed 모두 마지막 0.1M 학습 배송 $\ge 1.0$/episode \\
4. 배송 분산 & 마지막 0.1M에서 최다 agent 비율 $\le 0.7$ (배송 20건 미만이면 보류) \\
5. TOGGLE 남발 없음 & 요청되지 않은 선반 pickup이 같은 seed nopen의 50\% 이하 \\
\bottomrule
\end{tabularx}
\end{center}

%% ------------------------------------------------------------------
\section{초판 수정 범위}

10월 21일 발표본은 전면 수정하지 않는다. 서론·3장·5장은 구도를 정한 뒤 한 번에 다시 쓴다. 지금 고칠 것은 다음과 같다.
\begin{itemize}
  \item 9절의 ``E1-F 진행 중, 10월 24일 종료''를 ``10월 19일 1.536M에서 중단''으로 바로잡는다.
  \item 관련 연구에 Du et al.을 가장 가까운 연구로 넣고, Peer Learning·IDRL을 추가한다.
  \item 명제 3을 정책·가치·보상 예측 모두의 비식별 조건으로 넓힌다.
  \item 전제를 가설로 낮추고 명제 1·2를 부록으로 옮긴다.
  \item 깨진 ``부록??'' 참조와 \texttt{textbfRandom-gated} 오타를 고친다.
  \item 7장의 단일 seed 결과에서 성능 우위처럼 읽히는 문장을 줄인다.
\end{itemize}
6--7장 실험과 8--9장(수정 방향, 제출 계획)은 결과가 나온 뒤 고친다. 발표본 원고는 Overleaf에서 직접 고친 버전이라, 반영하려면 그 소스가 필요하다.

%% ------------------------------------------------------------------
\section{결정이 필요한 것}
\begin{itemize}
  \item pen005 판정을 보고 환경을 확정한다(10월 23일).
  \item 논문 구도를 ``식별 가능성 조건 + 조건별 probe''로 갈지 정한다.
  \item 메인 비교에 쓸 칸 두 개를 고른다(권장: 모서리·팀, 모서리·개인).
  \item 메인 학습 길이를 정한다(지금 예상 3M).
  \item 고정 $\tau$ 대신 상대 순위 판정을 검증할지 정한다.
\end{itemize}
열린 질문은 두 가지다. 인접 goal 칸에도 감점 .005가 통하는가. 개인 보상 칸에서 Oracle-team이 All-peer보다 낫기는 한가.

%% ------------------------------------------------------------------
\begin{thebibliography}{99}
\bibitem{agarwal2021} R. Agarwal, M. Schwarzer, P. S. Castro, A. Courville, M. G. Bellemare. Deep reinforcement learning at the edge of the statistical precipice. NeurIPS, 2021.
\bibitem{christianos2020} F. Christianos, L. Sch\"afer, S. V. Albrecht. Shared experience actor-critic for multi-agent reinforcement learning. NeurIPS, 2020.
\bibitem{christianos2021} F. Christianos, G. Papoudakis, A. Rahman, S. V. Albrecht. Scaling multi-agent reinforcement learning with selective parameter sharing. ICML, 2021.
\bibitem{derstroff2024} C. Derstroff, M. Cerrato, J. Brugger, J. Peters, S. Kramer. Peer learning: Learning complex policies in groups from scratch via action recommendations. AAAI, 2024.
\bibitem{du2025} H. Du, H. Nguyen, S. Thudumu, R. Vasa, K. Mouzakis. Goal-oriented multi-agent reinforcement learning for decentralized agent teams. arXiv:2511.11992, 2025.
\bibitem{hu2024} T. Hu, Z. Pu, X. Ai, T. Qiu, J. Yi. Measuring policy distance for multi-agent reinforcement learning. AAMAS, 2024.
\bibitem{idrl} Classifying ambiguous identities in hidden-role stochastic games with multi-agent reinforcement learning (IDRL). arXiv:2210.12896.
\bibitem{papoudakis2021} G. Papoudakis, F. Christianos, L. Sch\"afer, S. V. Albrecht. Benchmarking multi-agent deep reinforcement learning algorithms in cooperative tasks. NeurIPS Datasets and Benchmarks, 2021.
\bibitem{rusu2016} A. A. Rusu et al. Policy distillation. ICLR, 2016.
\bibitem{schulman2017} J. Schulman, F. Wolski, P. Dhariwal, A. Radford, O. Klimov. Proximal policy optimization algorithms. arXiv:1707.06347, 2017.
\bibitem{zhang2018} K. Zhang, Z. Yang, H. Liu, T. Zhang, T. Ba\c{s}ar. Fully decentralized multi-agent reinforcement learning with networked agents. ICML, 2018.
\end{thebibliography}

\appendix
\section{기록 번호와 파일}
수치의 근거는 저장소의 연구 기록(\texttt{EXPERIMENTS.md}, \texttt{DECISIONS.md})과 분석 결과 파일이다.
\begin{center}
\begin{tabularx}{\textwidth}{>{\raggedright\arraybackslash}p{2.3cm} L L}
\toprule
기록 & 내용 & 파일 \\
\midrule
E51, D74 & 2$\times$2 새 칸, 대기열, 판정 스크립트; 기반 통일 & \texttt{latent\_team/env\_candidates.py}, \texttt{run\_queue.ps1}, \texttt{pilot\_verdict.py} \\
E50, D73 & 식별 시점; 초판 수정 범위 & \texttt{output/reward\_coupling\_20261022/identification\_time.json} \\
E49, D72 & 보상 결합 대 정책 probe & \texttt{output/reward\_coupling\_20261022/results.json} \\
D71 & 독창성 재점검, 식별 가능성 구도 & \texttt{DECISIONS.md} \\
E46--E48, D70 & 감점 조정 & \texttt{output/nopen\_pilot\_20261022/}, \texttt{output/penalty\_sweep\_20261022/} \\
E43, E44, D68, D69 & 후보 pilot, pickup 정체 원인 & \texttt{output/pickup\_stall\_20261021/} \\
E36--E39 & 보상 probe $\hat r(o,a)$ & \texttt{output/reward\_probe\_20261019/} \\
E35 & Random-gated 구현 & \texttt{latent\_team/random\_gated.py} \\
D41, D45, D54 & 연구 목적, 지표 신뢰성 & \texttt{DECISIONS.md} \\
\bottomrule
\end{tabularx}
\end{center}

\end{document}
